% Endliche Faserzählung ohne Summen- oder Möbiustheorie.
\chapter{Endliche Fasern und kumulative Mächtigkeiten}

\Needspace{14\baselineskip}
\FormulaThmDeltaK[Gleiche endliche Kardinalzahlen liefern Bijektionen]{%
 \Finite A\dsep\Finite B\dsep\FiniteCard(A)=\FiniteCard(B)
 \vdash\exists f\,(f\colon A\bij B)
}{FiniteCardEqualityBijection}{%
 \DeltaRow{Mengen}{A\dsep B}\DeltaRow{Funktionen}{f}
}
\begin{tabproof}
 \proofstep{1}{\Finite A}{\rA}
 \proofstep{2}{\Finite B}{\rA}
 \proofstep{3}{\FiniteCard(A)=\FiniteCard(B)}{\rA}
 \proofstep{1}{\NatSeg{\FiniteCard(A)}\EqCard A}{\FormulaRefAuto{FiniteCardCounts}[thm]{1}}
 \proofstep{2}{\NatSeg{\FiniteCard(B)}\EqCard B}{\FormulaRefAuto{FiniteCardCounts}[thm]{2}}
 \proofstep{1}{A\EqCard\NatSeg{\FiniteCard(A)}}{\FormulaRefAuto{EqCardSymmetric}[thm]{4}}
 \proofstep{2,3}{\NatSeg{\FiniteCard(A)}\EqCard B}{\rIE{\FormulaRefAuto{a=b\vdash b=a}[thm]{3},5}}
 \proofstep{1,2,3}{A\EqCard B}{\FormulaRefAuto{EqCardTransitive}[thm]{6,7}}
 \proofstep{1,2,3}{\exists f\,(f\colon A\bij B)}{\FormulaRefAuto{Gleichmächtigkeit}[def]{8}}
\end{tabproof}

\Needspace{14\baselineskip}
\FormulaThmDeltaK[Bijektionen erhalten endliche Kardinalzahlen]{%
 \Finite A\dsep\Finite B\dsep\exists f\,(f\colon A\bij B)
 \vdash\FiniteCard(A)=\FiniteCard(B)
}{FiniteBijectionCardEquality}{%
 \DeltaRow{Mengen}{A\dsep B}\DeltaRow{Funktionen}{f}
}
\begin{tabproof}
 \proofstep{1}{\Finite A}{\rA}
 \proofstep{2}{\Finite B}{\rA}
 \proofstep{3}{\exists f\,(f\colon A\bij B)}{\rA}
 \proofstep{3}{A\EqCard B}{\FormulaRefAuto{Gleichmächtigkeit}[def]{3}}
 \proofstep{1}{\NatSeg{\FiniteCard(A)}\EqCard A}{\FormulaRefAuto{FiniteCardCounts}[thm]{1}}
 \proofstep{2}{\NatSeg{\FiniteCard(B)}\EqCard B}{\FormulaRefAuto{FiniteCardCounts}[thm]{2}}
 \proofstep{1,3}{\NatSeg{\FiniteCard(A)}\EqCard B}{\FormulaRefAuto{EqCardTransitive}[thm]{5,4}}
 \proofstep{1}{\FiniteCard(A)\in\mathbb N}{\FormulaRefAuto{FiniteCardNatural}[thm]{1}}
 \proofstep{2}{\FiniteCard(B)\in\mathbb N}{\FormulaRefAuto{FiniteCardNatural}[thm]{2}}
 \proofstep{1,2,3}{\FiniteCard(A)=\FiniteCard(B)}{\FormulaRefAuto{FiniteCardWitnessUnique}[thm]{8,9,7,6}}
\end{tabproof}

\Needspace{14\baselineskip}
\FormulaThmDeltaK[Endlichkeit der nichtleeren Potenzmenge reflektiert Endlichkeit]{%
 \Finite{\PowNE A}\vdash\Finite A
}{FiniteNonemptyPowerSetReflectsFiniteness}{%
 \DeltaRow{Mengen}{A\dsep x\dsep y}\DeltaRow{Funktionen}{j}
}
\begin{tabproofwide}
 \proofstepwidestar[1]{\Finite{\PowNE A}}{\rA}
 \proofstepwidestar[2]{x\in A}{\rA}
 \proofstepwidestar[2]{\{x\}\subseteq A}{\FormulaRefAuto{a\in A\vdash\{a\}\subseteq A}[thm]{2}}
 \proofstepwidestar[]{\{x\}\neq\varnothing}{\FormulaRefAuto{a\in A\vdash A\neq\varnothing}[thm]{\FormulaRefAuto{a\in\{a\}}[thm]}}
 \proofstepwidestar[2]{\{x\}\in\PowNE A}{\FormulaRefAuto{NonemptyPowersetMembership}[thm]{\rAI{3,4}}}
 \proofstepwidestar[]{\forall x\in A\,\{x\}\in\PowNE A}{\rUI{\rRI{2,5}}}
 \proofstepwidestar[]{\exists j\,(j\colon A\to\PowNE A\land\forall x\in A\,j(x)=\{x\})}{\FormulaRefAuto{TypedFunctionDefinitionPrinciple}[thm]{6}}
 \proofstepwidestar[8]{j\colon A\to\PowNE A\land\forall x\in A\,j(x)=\{x\}}{\rA}
 \proofstepwidestar[9]{y\in A\land j(x)=j(y)}{\rA}
 \proofstepwidestar[2,8]{j(x)=\{x\}}{\rRE{\rUE{\rAEb{8}},2}}
 \proofstepwidestar[8,9]{j(y)=\{y\}}{\rRE{\rUE{\rAEb{8}},\rAEa{9}}}
 \proofstepwidestar[2,8,9]{\{x\}=\{y\}}{\rIE{10,11,\rAEb{9}}}
 \proofstepwidestar[2,8,9]{x=y}{\FormulaRefAuto{\{a\}=\{b\}\eqvdash a=b}[thm]{12}}
 \proofstepwidestar[8]{\forall x,y\in A\,(j(x)=j(y)\rightarrow x=y)}{\rUI{\rRI{2,\rUI{\FormulaRefAuto{(P\land Q)\rightarrow R\eqvdash P\rightarrow(Q\rightarrow R)}[thm]{\rRI{9,13}}}}}}
 \proofstepwidestar[8]{j\colon A\inj\PowNE A}{\FormulaRefAuto{Injektive Funktion}[def]{\rAEa{8},14}}
 \proofstepwidestar[8]{j\colon A\bij j[A]}{\FormulaRefAuto{F\colon A\inj B\vdash F\colon A\bij F[A]}[thm]{15}}
 \proofstepwidestar[8]{j[A]\subseteq\PowNE A}{\FormulaRefAuto{F\colon A\to B\vdash F[A]\subseteq B}[thm]{\rAEa{8}}}
 \proofstepwidestar[1,8]{\Finite{j[A]}}{\FormulaRefAuto{FiniteSubset}[thm]{17,1}}
 \proofstepwidestar[8]{j^{-1}\colon j[A]\bij A}{\FormulaRefAuto{InverseFunctionBijection}[thm]{16}}
 \proofstepwidestar[1,8]{\Finite A}{\FormulaRefAuto{FiniteBijectionTransport}[thm]{18,19}}
 \proofstepwidestar[1]{\Finite A}{\rEE{7,8,20}}
\end{tabproofwide}

\Needspace{14\baselineskip}
\FormulaThmDeltaK[Fortsetzung einer endlichen partiellen Injektion]{%
 \begin{aligned}[t]
 &\Finite U\dsep U\subseteq A\dsep u\colon U\inj B\dsep
 \exists v\,(v\colon A\bij B)\vdash{}\\[-2pt]
 &\exists v\,(v\colon A\bij B\land\forall x\in U\,v(x)=u(x))
 \end{aligned}
}{FinitePartialBijectionExtension}{%
 \DeltaRow{Mengen}{U\dsep A\dsep B\dsep C\dsep a\dsep x}
 \DeltaRow{Funktionen}{u\dsep v\dsep w}
}
Für den Induktionsbeweis steht
\(P(C)\) für
\(C\subseteq U\rightarrow
 \exists v\,(v\colon A\bij B\land\forall x\in C\,v(x)=u(x))\).
\begin{tabproofsplitwide}
 \proofpartwide{Leerer Anfang}
 \proofstepwidestar[1]{U\subseteq A}{\rA}
 \proofstepwidestar[2]{u\colon U\inj B}{\rA}
 \proofstepwidestar[3]{\exists v\,(v\colon A\bij B)}{\rA}
 \proofstepwidestar[4]{v\colon A\bij B}{\rA}
 \proofstepwidestar[5]{x\in\varnothing}{\rA}
 \proofstepwidestar[]{x\notin\varnothing}{\FormulaRefAuto{x\not\in\varnothing}[thm]}
 \proofstepwidestar[5]{\bot}{\rBI{6,5}}
 \proofstepwidestar[5]{v(x)=u(x)}{\FormulaRefAuto{\bot\vdash P}[thm]{7}}
 \proofstepwidestar[]{\forall x\in\varnothing\,v(x)=u(x)}{\rUI{\rRI{5,8}}}
 \proofstepwidestar[4]{\exists v\,(v\colon A\bij B\land\forall x\in\varnothing\,v(x)=u(x))}{\rEI{\rAI{4,9}}}
 \proofstepwidestar[3]{\exists v\,(v\colon A\bij B\land\forall x\in\varnothing\,v(x)=u(x))}{\rEE{3,4,10}}
 \proofstepwidestar[12]{\varnothing\subseteq U}{\rA}
 \proofstepwidestar[3]{P(\varnothing)}{\rRI{12,11}}
 \closeproofpartwide
 \Needspace{5\baselineskip}
 \proofpartwide{Adjunktion}
 \setcounter{proofstepnr}{13}
 \proofstepwidestar[14]{\Finite C\land P(C)\land a\notin C}{\rA}
 \proofstepwidestar[15]{C\cup\{a\}\subseteq U}{\rA}
 \proofstepwidestar[]{C\subseteq C\cup\{a\}}{\FormulaRefAuto{A\subseteq A\cup B}[thm]}
 \proofstepwidestar[15]{C\subseteq U}{\FormulaRefAuto{A\subseteq B\dsep B\subseteq C\vdash A\subseteq C}[thm]{16,15}}
 \proofstepwidestar[]{a\in C\cup\{a\}}{\FormulaRefAuto{a\in A\cup\{a\}}[thm]}
 \proofstepwidestar[15]{a\in U}{\FormulaRefAuto{A\subseteq B\dsep x\in A\vdash x\in B}[thm]{15,18}}
 \proofstepwidestar[14,15]{\exists v\,(v\colon A\bij B\land\forall x\in C\,v(x)=u(x))}{\rRE{\rAEa{\rAEb{14}},17}}
 \proofstepwidestar[21]{v\colon A\bij B\land\forall x\in C\,v(x)=u(x)}{\rA}
 \proofstepwidestar[1,2,14,15,21]{\exists w\,(w\colon A\bij B\land\forall x\in C\cup\{a\}\,w(x)=u(x))}{\FormulaRefAuto{BijectionAgreementExtensionStep}[thm]{1,17,19,\rAEb{\rAEb{14}},2,\rAEa{21},\rAEb{21}}}
 \proofstepwidestar[1,2,14,15]{\exists w\,(w\colon A\bij B\land\forall x\in C\cup\{a\}\,w(x)=u(x))}{\rEE{20,21,22}}
 \proofstepwidestar[1,2,14]{P(C\cup\{a\})}{\rRI{15,23}}
 \closeproofpartwide
 \proofpartwide{Induktionsschluss}
 \setcounter{proofstepnr}{24}
 \proofstepwidestar[25]{\Finite U}{\rA}
 \proofstepwidestar[1,2,3,25]{P(U)}{\FormulaRefAuto{FiniteInductionRule}[ax]{13,[14]\vdots24,25}}
 \proofstepwidestar[]{U\subseteq U}{\FormulaRefAuto{A\subseteq A}[thm]}
 \proofstepwidestar[1,2,3,25]{\exists v\,(v\colon A\bij B\land\forall x\in U\,v(x)=u(x))}{\rRE{26,27}}
 \closeproofpartwide
\end{tabproofsplitwide}

\Needspace{14\baselineskip}
\FormulaThmDeltaK[Endliche gleichmächtige Teile können gekürzt werden]{%
 \begin{aligned}[t]
 &\Finite C\dsep C\subseteq A\dsep D\subseteq B\dsep
 \exists u\,(u\colon C\bij D)\dsep\exists v\,(v\colon A\bij B)\vdash{}\\[-2pt]
 &\exists h\,(h\colon A\setminus C\bij B\setminus D)
 \end{aligned}
}{FiniteComplementBijection}{%
 \DeltaRow{Mengen}{A\dsep B\dsep C\dsep D\dsep x\dsep z}
 \DeltaRow{Funktionen}{u\dsep v\dsep h}
}
\begin{tabproofsplitwide}
 \proofpartwide{Fortsetzung der Teilbijektion}
 \proofstepwidestar[1]{\Finite C\land C\subseteq A\land D\subseteq B}{\rA}
 \proofstepwidestar[2]{\exists u\,(u\colon C\bij D)}{\rA}
 \proofstepwidestar[3]{\exists v\,(v\colon A\bij B)}{\rA}
 \proofstepwidestar[4]{u\colon C\bij D}{\rA}
 \proofstepwidestar[1,4]{u\colon C\to B}{\FormulaRefAuto{F\colon A\to B\dsep B\subseteq C\vdash F\colon A\to C}[thm]{4,\rAEb{\rAEb{1}}}}
 \proofstepwidestar[4]{\forall x,z\in C\,(u(x)=u(z)\rightarrow x=z)}{\FormulaRefAuto{Bijektive Funktion}[def]{4}}
 \proofstepwidestar[1,4]{u\colon C\inj B}{\FormulaRefAuto{Injektive Funktion}[def]{5,6}}
 \proofstepwidestar[1,3,4]{\exists v\,(v\colon A\bij B\land\forall x\in C\,v(x)=u(x))}{\FormulaRefAuto{FinitePartialBijectionExtension}[thm]{\rAEa{1},\rAEa{\rAEb{1}},7,3}}
 \proofstepwidestar[9]{v\colon A\bij B\land\forall x\in C\,v(x)=u(x)}{\rA}
 \proofstepwidestar[10]{x\in A}{\rA}
 \proofstepwidestar[11]{x\in C}{\rA}
 \proofstepwidestar[4,11]{u(x)\in D}{\FormulaRefAuto{F\colon A\bij B\dsep x\in A\vdash F(x)\in B}[thm]{4,11}}
 \proofstepwidestar[9,11]{v(x)=u(x)}{\rRE{\rUE{\rAEb{9}},11}}
 \proofstepwidestar[4,9,11]{v(x)\in D}{\rIE{\FormulaRefAuto{a=b\vdash b=a}[thm]{13},12}}
 \proofstepwidestar[4,9]{x\in C\rightarrow v(x)\in D}{\rRI{11,14}}
 \proofstepwidestar[16]{v(x)\in D}{\rA}
 \proofstepwidestar[4,16]{\exists z\in C\,u(z)=v(x)}{\FormulaRefAuto{y\in B\vdash\exists x\in A\;F(x)=y}[ax]{4,16}}
 \proofstepwidestar[18]{z\in C\land u(z)=v(x)}{\rA}
 \proofstepwidestar[1,18]{z\in A}{\FormulaRefAuto{A\subseteq B\dsep x\in A\vdash x\in B}[thm]{\rAEa{\rAEb{1}},\rAEa{18}}}
 \proofstepwidestar[9,18]{v(z)=u(z)}{\rRE{\rUE{\rAEb{9}},\rAEa{18}}}
 \proofstepwidestar[9,18]{v(z)=v(x)}{\rIE{\rAEb{18},20}}
 \proofstepwidestar[1,9,10,18]{z=x}{\FormulaRefAuto{F\colon A\inj B\dsep x\in A\dsep y\in A\dsep F(x)=F(y)\vdash x=y}[ax]{\rAEa{9},19,10,21}}
 \proofstepwidestar[1,9,10,18]{x\in C}{\rIE{22,\rAEa{18}}}
 \proofstepwidestar[1,4,9,10,16]{x\in C}{\rEE{17,18,23}}
 \proofstepwidestar[1,4,9,10]{v(x)\in D\rightarrow x\in C}{\rRI{16,24}}
 \proofstepwidestar[1,4,9,10]{x\in C\leftrightarrow v(x)\in D}{\rLRI{15,25}}
 \closeproofpartwide
 \proofpartwide{Einschränkung auf die Komplemente}
 \setcounter{proofstepnr}{26}
 \proofstepwidestar[9,10]{v(x)\in B}{\FormulaRefAuto{F\colon A\bij B\dsep x\in A\vdash F(x)\in B}[thm]{\rAEa{9},10}}
 \proofstepwidestar[1,4,9,10]{x\notin C\leftrightarrow v(x)\notin D}{\FormulaRefAuto{P\leftrightarrow Q\eqvdash\neg P\leftrightarrow\neg Q}[thm]{26}}
 \proofstepwidestar[10]{x\in A\setminus C\leftrightarrow x\notin C}{\FormulaRefAuto{BoundedDifferenceMembership}[thm]{10}}
 \proofstepwidestar[9,10]{v(x)\in B\setminus D\leftrightarrow v(x)\notin D}{\FormulaRefAuto{BoundedDifferenceMembership}[thm]{27}}
 \proofstepwidestar[1,4,9,10]{x\in A\setminus C\leftrightarrow v(x)\in B\setminus D}{\rChain{29,28,\FormulaRefAuto{P\leftrightarrow Q\vdash Q\leftrightarrow P}[thm]{30}}}
 \proofstepwidestar[1,4,9]{\forall x\in A\,(x\in A\setminus C\leftrightarrow v(x)\in B\setminus D)}{\rUI{\rRI{10,31}}}
 \proofstepwidestar[]{A\setminus C\subseteq A\land B\setminus D\subseteq B}{\rAI{\FormulaRefAuto{A\setminus B\subseteq A}[thm],\FormulaRefAuto{A\setminus B\subseteq A}[thm]}}
 \proofstepwidestar[1,4,9]{v\restriction_{A\setminus C}\colon A\setminus C\bij B\setminus D}{\FormulaRefAuto{PointwiseCompatibilityImpliesBijectiveRestriction}[thm]{\rAEa{9},\rAEa{33},\rAEb{33},32}}
 \proofstepwidestar[1,4,9]{\exists h\,(h\colon A\setminus C\bij B\setminus D)}{\rEI{34}}
 \proofstepwidestar[1,3,4]{\exists h\,(h\colon A\setminus C\bij B\setminus D)}{\rEE{8,9,35}}
 \proofstepwidestar[1,2,3]{\exists h\,(h\colon A\setminus C\bij B\setminus D)}{\rEE{2,4,36}}
 \closeproofpartwide
\end{tabproofsplitwide}
\Needspace{14\baselineskip}
\FormulaThmDeltaK[Induktion über eine endliche partielle Ordnung]{%
 \begin{aligned}[t]
 &\PartOrd{Q,\leq}\dsep\Finite Q\dsep{}\\[-2pt]
 &\forall p\in Q\,\Bigl(
   (\forall q\in Q\,(q\leq p\land q\neq p\rightarrow P(q)))
    \rightarrow P(p)\Bigr)\vdash\forall p\in Q\,P(p)
 \end{aligned}
}{FinitePosetInduction}{%
 \DeltaRow{Mengen}{Q\dsep N\dsep p\dsep q}\DeltaRow{Relationsoperatoren}{\leq}
 \DeltaRow{Prädikat}{P}
}
Setze \(N=\{p\in Q\mid\neg P(p)\}\).
\begin{tabproofwide}
 \proofstepwidestar[1]{\PartOrd{Q,\leq}}{\rA}
 \proofstepwidestar[2]{\Finite Q}{\rA}
 \proofstepwidestar[3]{\forall p\in Q\,((\forall q\in Q\,(q\leq p\land q\neq p\rightarrow P(q)))\rightarrow P(p))}{\rA}
 \proofstepwidestar[]{N\subseteq Q}{\FormulaRefAuto{\{x\in A\mid P(x)\}\subseteq A}[thm]}
 \proofstepwidestar[2]{\Finite N}{\FormulaRefAuto{FiniteSubset}[thm]{4,2}}
 \proofstepwidestar[6]{N\neq\varnothing}{\rA}
 \proofstepwidestar[1,2,6]{\exists p\in N\,\forall q\in N\,(q\leq p\rightarrow q=p)}{\FormulaRefAuto{FinitePosetMinimalElement}[thm]{1,5,6,4}}
 \proofstepwidestar[8]{p\in N\land\forall q\in N\,(q\leq p\rightarrow q=p)}{\rA}
 \proofstepwidestar[8]{p\in Q\land\neg P(p)}{\FormulaRefAuto{x\in\{u\in A\mid P(u)\}\eqvdash x\in A\land P(x)}[thm]{\rAEa{8}}}
 \proofstepwidestar[10]{q\in Q\land q\leq p\land q\neq p}{\rA}
 \proofstepwidestar[11]{\neg P(q)}{\rA}
 \proofstepwidestar[10,11]{q\in N}{\FormulaRefAuto{x\in A\dsep P(x)\vdash x\in\{x\in A\mid P(x)\}}[thm]{\rAEa{10},11}}
 \proofstepwidestar[8,10,11]{q=p}{\rRE{\rRE{\rUE{\rAEb{8}},12},\rAEa{\rAEb{10}}}}
 \proofstepwidestar[8,10,11]{\bot}{\rBI{\rAEb{\rAEb{10}},13}}
 \proofstepwidestar[8,10]{P(q)}{\rCE{11,14}}
 \proofstepwidestar[8]{\forall q\in Q\,(q\leq p\land q\neq p\rightarrow P(q))}{\rUI{\FormulaRefAuto{(P\land Q)\rightarrow R\eqvdash P\rightarrow(Q\rightarrow R)}[thm]{\rRI{10,15}}}}
 \proofstepwidestar[3,8]{(\forall q\in Q\,(q\leq p\land q\neq p\rightarrow P(q)))\rightarrow P(p)}{\rRE{\rUE{3},\rAEa{9}}}
 \proofstepwidestar[3,8]{P(p)}{\rRE{17,16}}
 \proofstepwidestar[3,8]{\bot}{\rBI{\rAEb{9},18}}
 \proofstepwidestar[1,2,3,6]{\bot}{\rEE{7,8,19}}
 \proofstepwidestar[1,2,3]{N=\varnothing}{\rCE{6,20}}
 \proofstepwidestar[22]{p\in Q}{\rA}
 \proofstepwidestar[23]{\neg P(p)}{\rA}
 \proofstepwidestar[22,23]{p\in N}{\FormulaRefAuto{x\in A\dsep P(x)\vdash x\in\{x\in A\mid P(x)\}}[thm]{22,23}}
 \proofstepwidestar[1,2,3,22,23]{p\in\varnothing}{\rIE{21,24}}
 \proofstepwidestar[]{p\notin\varnothing}{\FormulaRefAuto{x\not\in\varnothing}[thm]}
 \proofstepwidestar[1,2,3,22,23]{\bot}{\rBI{26,25}}
 \proofstepwidestar[1,2,3,22]{P(p)}{\rCE{23,27}}
 \proofstepwidestar[1,2,3]{\forall p\in Q\,P(p)}{\rUI{\rRI{22,28}}}
\end{tabproofwide}

Für eine partielle Ordnung auf \(Q\) schreiben wir
\[
 I_p=\{q\in Q\mid q\leq p\},\qquad
 J_p=\{q\in Q\mid q\leq p\land q\neq p\}.
\]
Dies sind die untere Menge einschließlich beziehungsweise ausschließlich
des Punktes \(p\).

\Needspace{14\baselineskip}
\FormulaThmDeltaK[Eine untere Faserfamilie zerfällt am Endpunkt]{%
 \begin{aligned}[t]
 &\PartOrd{Q,\leq}\dsep r\colon S\to Q\dsep p\in Q\vdash{}\\[-2pt]
 &r^{-1}[J_p]\subseteq r^{-1}[I_p]\land
 \bigl(r^{-1}[I_p]\setminus r^{-1}[J_p]\bigr)=r^{-1}[\{p\}]
 \end{aligned}
}{PrincipalLowerFiberDecomposition}{%
 \DeltaRow{Mengen}{Q\dsep S\dsep p\dsep q\dsep x\dsep I_p\dsep J_p}
 \DeltaRow{Relationsoperatoren}{\leq}\DeltaRow{Funktionen}{r}
}
\begin{tabproofsplitwide}
 \proofpartwide{Die beiden unteren Mengen}
 \proofstepwidestar[1]{\PartOrd{Q,\leq}}{\rA}
 \proofstepwidestar[2]{r\colon S\to Q}{\rA}
 \proofstepwidestar[3]{p\in Q}{\rA}
 \proofstepwidestar[]{I_p\subseteq Q\land J_p\subseteq Q}{\rAI{\FormulaRefAuto{\{x\in A\mid P(x)\}\subseteq A}[thm],\FormulaRefAuto{\{x\in A\mid P(x)\}\subseteq A}[thm]}}
 \proofstepwidestar[5]{q\in J_p}{\rA}
 \proofstepwidestar[5]{q\in Q\land q\leq p\land q\neq p}{\FormulaRefAuto{x\in\{u\in A\mid P(u)\}\eqvdash x\in A\land P(x)}[thm]{5}}
 \proofstepwidestar[5]{q\in I_p}{\FormulaRefAuto{x\in A\dsep P(x)\vdash x\in\{x\in A\mid P(x)\}}[thm]{\rAEa{6},\rAEa{\rAEb{6}}}}
 \proofstepwidestar[]{J_p\subseteq I_p}{\FormulaRefAuto{SubsetIntroductionRule}[thm]{[5]\vdots7}}
 \proofstepwidestar[2]{r^{-1}[J_p]\subseteq r^{-1}[I_p]}{\FormulaRefAuto{F\colon A\to B\dsep C\subseteq D\dsep D\subseteq B\vdash F^{-1}[C]\subseteq F^{-1}[D]}[thm]{2,8,\rAEa{4}}}
 \closeproofpartwide
 \proofpartwide{Die verbleibenden Elemente liegen in der Endpunktfaser}
 \setcounter{proofstepnr}{9}
 \proofstepwidestar[10]{x\in r^{-1}[I_p]\setminus r^{-1}[J_p]}{\rA}
 \proofstepwidestar[10]{x\in r^{-1}[I_p]\land x\notin r^{-1}[J_p]}{\FormulaRefAuto{x\in A\setminus B\eqvdash x\in A\land x\notin B}[thm]{10}}
 \proofstepwidestar[2,10]{x\in S\land r(x)\in I_p}{\rRE{\rLREa{\FormulaRefAuto{F\colon A\to B\dsep C\subseteq B\vdash\bigl(x\in F^{-1}[C]\leftrightarrow(x\in A\land F(x)\in C)\bigr)}[thm]{2,\rAEa{4}}},\rAEa{11}}}
 \proofstepwidestar[2,10]{r(x)\in Q\land r(x)\leq p}{\FormulaRefAuto{x\in\{u\in A\mid P(u)\}\eqvdash x\in A\land P(x)}[thm]{\rAEb{12}}}
 \proofstepwidestar[14]{r(x)\neq p}{\rA}
 \proofstepwidestar[2,10,14]{r(x)\in J_p}{\FormulaRefAuto{x\in A\dsep P(x)\vdash x\in\{x\in A\mid P(x)\}}[thm]{\rAEa{13},\rAI{\rAEb{13},14}}}
 \proofstepwidestar[2,10,14]{x\in r^{-1}[J_p]}{\rRE{\rLREb{\FormulaRefAuto{F\colon A\to B\dsep C\subseteq B\vdash\bigl(x\in F^{-1}[C]\leftrightarrow(x\in A\land F(x)\in C)\bigr)}[thm]{2,\rAEb{4}}},\rAI{\rAEa{12},15}}}
 \proofstepwidestar[2,10,14]{\bot}{\rBI{\rAEb{11},16}}
 \proofstepwidestar[2,10]{r(x)=p}{\rCE{14,17}}
 \proofstepwidestar[2,3,10]{x\in r^{-1}[\{p\}]}{\rRE{\rLREb{\FormulaRefAuto{FunctionFiberMembership}[thm]{2,3}},\rAI{\rAEa{12},18}}}
 \proofstepwidestar[2,3]{r^{-1}[I_p]\setminus r^{-1}[J_p]\subseteq r^{-1}[\{p\}]}{\FormulaRefAuto{SubsetIntroductionRule}[thm]{[10]\vdots19}}
 \closeproofpartwide
 \proofpartwide{Die Endpunktfaser bleibt vollständig übrig}
 \setcounter{proofstepnr}{20}
 \proofstepwidestar[21]{x\in r^{-1}[\{p\}]}{\rA}
 \proofstepwidestar[2,3,21]{x\in S\land r(x)=p}{\rRE{\rLREa{\FormulaRefAuto{FunctionFiberMembership}[thm]{2,3}},21}}
 \proofstepwidestar[1,3]{p\leq p}{\FormulaRefAuto{EqRelReflexivity}[ax]{1,3}}
 \proofstepwidestar[1,3]{p\in I_p}{\FormulaRefAuto{x\in A\dsep P(x)\vdash x\in\{x\in A\mid P(x)\}}[thm]{3,23}}
 \proofstepwidestar[1,2,3,21]{r(x)\in I_p}{\rIE{\FormulaRefAuto{a=b\vdash b=a}[thm]{\rAEb{22}},24}}
 \proofstepwidestar[1,2,3,21]{x\in r^{-1}[I_p]}{\rRE{\rLREb{\FormulaRefAuto{F\colon A\to B\dsep C\subseteq B\vdash\bigl(x\in F^{-1}[C]\leftrightarrow(x\in A\land F(x)\in C)\bigr)}[thm]{2,\rAEa{4}}},\rAI{\rAEa{22},25}}}
 \proofstepwidestar[27]{x\in r^{-1}[J_p]}{\rA}
 \proofstepwidestar[2,27]{x\in S\land r(x)\in J_p}{\rRE{\rLREa{\FormulaRefAuto{F\colon A\to B\dsep C\subseteq B\vdash\bigl(x\in F^{-1}[C]\leftrightarrow(x\in A\land F(x)\in C)\bigr)}[thm]{2,\rAEb{4}}},27}}
 \proofstepwidestar[2,27]{r(x)\in Q\land r(x)\leq p\land r(x)\neq p}{\FormulaRefAuto{x\in\{u\in A\mid P(u)\}\eqvdash x\in A\land P(x)}[thm]{\rAEb{28}}}
 \proofstepwidestar[2,3,21,27]{\bot}{\rBI{\rAEb{\rAEb{29}},\rAEb{22}}}
 \proofstepwidestar[2,3,21]{x\notin r^{-1}[J_p]}{\rCI{27,30}}
 \proofstepwidestar[1,2,3,21]{x\in r^{-1}[I_p]\setminus r^{-1}[J_p]}{\FormulaRefAuto{x\in A\setminus B\eqvdash x\in A\land x\notin B}[thm]{\rAI{26,31}}}
 \proofstepwidestar[1,2,3]{r^{-1}[\{p\}]\subseteq r^{-1}[I_p]\setminus r^{-1}[J_p]}{\FormulaRefAuto{SubsetIntroductionRule}[thm]{[21]\vdots32}}
 \proofstepwidestar[1,2,3]{r^{-1}[I_p]\setminus r^{-1}[J_p]=r^{-1}[\{p\}]}{\FormulaRefAuto{A\subseteq B\dsep B\subseteq A\vdash A=B}[thm]{20,33}}
 \proofstepwidestar[1,2,3]{\begin{aligned}[t]&r^{-1}[J_p]\subseteq r^{-1}[I_p]\land{}\\[-2pt]&(r^{-1}[I_p]\setminus r^{-1}[J_p])=r^{-1}[\{p\}]\end{aligned}}{\rAI{9,34}}
 \closeproofpartwide
\end{tabproofsplitwide}

\Needspace{14\baselineskip}
\FormulaThmDeltaK[Kumulative Fasergrößen bestimmen die einzelnen Fasern]{%
 \begin{aligned}[t]
 &\PartOrd{Q,\leq}\dsep\Finite Q\dsep\Finite S\dsep\Finite T\dsep
 r\colon S\to Q\dsep t\colon T\to Q\dsep{}\\[-2pt]
 &\forall p\in Q\,\FiniteCard(r^{-1}[I_p])=\FiniteCard(t^{-1}[I_p])\vdash{}\\[-2pt]
 &\forall p\in Q\,\exists u\,(u\colon r^{-1}[\{p\}]\bij t^{-1}[\{p\}])
 \end{aligned}
}{FinitePosetCumulativeFiberBijection}{%
 \DeltaRow{Mengen}{Q\dsep S\dsep T\dsep p\dsep q\dsep I_p\dsep J_p\dsep A_p\dsep B_p\dsep A^-_p\dsep B^-_p}
 \DeltaRow{Relationsoperatoren}{\leq}\DeltaRow{Funktionen}{r\dsep t\dsep u}
}
Im Beweis ist \(P(p)\) die Schlussformel
\(\exists u\,(u\colon r^{-1}[\{p\}]\bij t^{-1}[\{p\}])\).
Zur Abkürzung setzen wir für diese Tabelle
\[
 A_p=r^{-1}[I_p],\quad B_p=t^{-1}[I_p],\quad
 A^-_p=r^{-1}[J_p],\quad B^-_p=t^{-1}[J_p].
\]
\begin{tabproofwide}
 \proofstepwidestar[1]{\PartOrd{Q,\leq}\land\Finite Q}{\rA}
 \proofstepwidestar[2]{\Finite S\land\Finite T}{\rA}
 \proofstepwidestar[3]{r\colon S\to Q\land t\colon T\to Q}{\rA}
 \proofstepwidestar[4]{\forall p\in Q\,\FiniteCard(A_p)=\FiniteCard(B_p)}{\rA}
 \proofstepwidestar[5]{p\in Q}{\rA}
 \proofstepwidestar[6]{\forall q\in Q\,(q\leq p\land q\neq p\rightarrow P(q))}{\rA}
 \proofstepwidestar[]{I_p\subseteq Q\land J_p\subseteq Q}{\rAI{\FormulaRefAuto{\{x\in A\mid P(x)\}\subseteq A}[thm],\FormulaRefAuto{\{x\in A\mid P(x)\}\subseteq A}[thm]}}
 \proofstepwidestar[8]{q\in J_p}{\rA}
 \proofstepwidestar[8]{q\in Q\land q\leq p\land q\neq p}{\FormulaRefAuto{x\in\{u\in A\mid P(u)\}\eqvdash x\in A\land P(x)}[thm]{8}}
 \proofstepwidestar[6,8]{P(q)}{\rRE{\rRE{\rUE{6},\rAEa{9}},\rAEb{9}}}
 \proofstepwidestar[6]{\forall q\in J_p\,P(q)}{\rUI{\rRI{8,10}}}
 \proofstepwidestar[3,6]{\exists h\,(h\colon A^-_p\bij B^-_p)}{\FormulaRefAuto{FiberwiseRestrictedBijectionExists}[thm]{\rAEa{3},\rAEb{3},\rAEb{7},11}}
 \proofstepwidestar[3]{A_p\subseteq S\land B_p\subseteq T\land A^-_p\subseteq S}{\rAI{\FormulaRefAuto{F\colon A\to B\dsep C\subseteq B\vdash F^{-1}[C]\subseteq A}[thm]{\rAEa{3},\rAEa{7}},\rAI{\FormulaRefAuto{F\colon A\to B\dsep C\subseteq B\vdash F^{-1}[C]\subseteq A}[thm]{\rAEb{3},\rAEa{7}},\FormulaRefAuto{F\colon A\to B\dsep C\subseteq B\vdash F^{-1}[C]\subseteq A}[thm]{\rAEa{3},\rAEb{7}}}}}
 \proofstepwidestar[2,3]{\Finite{A_p}\land\Finite{B_p}\land\Finite{A^-_p}}{\rAI{\FormulaRefAuto{FiniteSubset}[thm]{\rAEa{13},\rAEa{2}},\rAI{\FormulaRefAuto{FiniteSubset}[thm]{\rAEa{\rAEb{13}},\rAEb{2}},\FormulaRefAuto{FiniteSubset}[thm]{\rAEb{\rAEb{13}},\rAEa{2}}}}}
 \proofstepwidestar[4,5]{\FiniteCard(A_p)=\FiniteCard(B_p)}{\rRE{\rUE{4},5}}
 \proofstepwidestar[2,3,4,5]{\exists v\,(v\colon A_p\bij B_p)}{\FormulaRefAuto{FiniteCardEqualityBijection}[thm]{\rAEa{14},\rAEa{\rAEb{14}},15}}
 \proofstepwidestar[1,3,5]{A^-_p\subseteq A_p\land(A_p\setminus A^-_p)=r^{-1}[\{p\}]}{\FormulaRefAuto{PrincipalLowerFiberDecomposition}[thm]{\rAEa{1},\rAEa{3},5}}
 \proofstepwidestar[1,3,5]{B^-_p\subseteq B_p\land(B_p\setminus B^-_p)=t^{-1}[\{p\}]}{\FormulaRefAuto{PrincipalLowerFiberDecomposition}[thm]{\rAEa{1},\rAEb{3},5}}
 \proofstepwidestar[1,2,3,4,5,6]{\exists u\,(u\colon A_p\setminus A^-_p\bij B_p\setminus B^-_p)}{\FormulaRefAuto{FiniteComplementBijection}[thm]{\rAEb{\rAEb{14}},\rAEa{17},\rAEa{18},12,16}}
 \proofstepwidestar[1,2,3,4,5,6]{P(p)}{\rIE{\rAEb{17},\rAEb{18},19}}
 \proofstepwidestar[1,2,3,4]{\forall p\in Q\,((\forall q\in Q\,(q\leq p\land q\neq p\rightarrow P(q)))\rightarrow P(p))}{\rUI{\rRI{5,\rRI{6,20}}}}
 \proofstepwidestar[1,2,3,4]{\forall p\in Q\,P(p)}{\FormulaRefAuto{FinitePosetInduction}[thm]{\rAEa{1},\rAEb{1},21}}
\end{tabproofwide}

\Needspace{14\baselineskip}
\FormulaThmDeltaK[Rücktransport unterer Mengen und einzelner Fasern]{%
 \begin{aligned}[t]
 &\psi\colon Q\bij R\dsep t\colon T\to R\dsep p\in Q\dsep{}\\[-2pt]
 &\forall a,b\in Q\,(a\leq b\leftrightarrow\psi(a)\preceq\psi(b))\vdash{}\\[-2pt]
 &v^{-1}[I_p]=t^{-1}[I'_{\psi(p)}]\land
 v^{-1}[\{p\}]=t^{-1}[\{\psi(p)\}]
 \end{aligned}
}{OrderPullbackLowerFibers}{%
 \DeltaRow{Mengen}{Q\dsep R\dsep T\dsep p\dsep x\dsep a\dsep b}
 \DeltaRow{Relationsoperatoren}{\leq\dsep\preceq}
 \DeltaRow{Funktionen}{\psi\dsep t\dsep v=\psi^{-1}\circ t}
 \DeltaRow{Abkürzungen}{I_p=\{a\in Q\mid a\leq p\}}
 \DeltaRow{}{I'_q=\{b\in R\mid b\preceq q\}}
}
\Needspace{7\baselineskip}
\begin{tabproofsplitwide}
 \proofpartwide{Punktweiser Rücktransport}
 \proofstepwidestar[1]{\psi\colon Q\bij R}{\rA}
 \proofstepwidestar[2]{t\colon T\to R}{\rA}
 \proofstepwidestar[3]{p\in Q}{\rA}
 \proofstepwidestar[4]{\forall a,b\in Q\,(a\leq b\leftrightarrow\psi(a)\preceq\psi(b))}{\rA}
 \proofstepwidestar[1]{\psi^{-1}\colon R\to Q}{\FormulaRefAuto{InverseFunctionType}[thm]{1}}
 \proofstepwidestar[1,2]{v\colon T\to Q}{\FormulaRefAuto{CompositionDef}[def]{2,5}}
 \proofstepwidestar[1,3]{\psi(p)\in R}{\FormulaRefAuto{F\colon A\bij B\dsep x\in A\vdash F(x)\in B}[thm]{1,3}}
 \proofstepwidestar[8]{x\in T}{\rA}
 \proofstepwidestar[2,8]{t(x)\in R}{\FormulaRefAuto{F\colon A\to B\dsep x\in A\vdash F(x)\in B}[thm]{2,8}}
 \proofstepwidestar[1,2,8]{v(x)\in Q}{\FormulaRefAuto{F\colon A\to B\dsep x\in A\vdash F(x)\in B}[thm]{6,8}}
 \proofstepwidestar[1,2,8]{v(x)=\psi^{-1}(t(x))}{\FormulaRefAuto{CompositionDef}[def]{2,5,8}}
 \proofstepwidestar[1,2,8]{\psi(\psi^{-1}(t(x)))=t(x)}{\FormulaRefAuto{InverseFunctionRightInverse}[thm]{1,9}}
 \proofstepwidestar[1,2,8]{\psi(v(x))=t(x)}{\rIE{\FormulaRefAuto{a=b\vdash b=a}[thm]{11},12}}
 \proofstepwidestar[1,2,3,4,8]{v(x)\leq p\leftrightarrow\psi(v(x))\preceq\psi(p)}{\rRE{\rUE{\rRE{\rUE{4},10}},3}}
 \proofstepwidestar[1,2,3,4,8]{v(x)\leq p\leftrightarrow t(x)\preceq\psi(p)}{\rIE{13,14}}
 \proofstepwidestar[16]{v(x)=p}{\rA}
 \proofstepwidestar[1,2,8,16]{t(x)=\psi(p)}{\rIE{16,\FormulaRefAuto{a=b\vdash b=a}[thm]{13}}}
 \proofstepwidestar[1,2,8]{v(x)=p\rightarrow t(x)=\psi(p)}{\rRI{16,17}}
 \proofstepwidestar[19]{t(x)=\psi(p)}{\rA}
 \proofstepwidestar[1,2,8,19]{\psi(v(x))=\psi(p)}{\rIE{19,13}}
 \proofstepwidestar[1,2,3,8,19]{v(x)=p}{\FormulaRefAuto{F\colon A\inj B\dsep x\in A\dsep y\in A\dsep F(x)=F(y)\vdash x=y}[ax]{1,10,3,20}}
 \proofstepwidestar[1,2,3,8]{t(x)=\psi(p)\rightarrow v(x)=p}{\rRI{19,21}}
 \proofstepwidestar[1,2,3,8]{v(x)=p\leftrightarrow t(x)=\psi(p)}{\rLRI{18,22}}
 \closeproofpartwide
 \proofpartwide{Gleichheit der Urbildmengen}
 \setcounter{proofstepnr}{23}
 \proofstepwidestar[]{I_p\subseteq Q\land I'_{\psi(p)}\subseteq R}{\rAI{\FormulaRefAuto{\{x\in A\mid P(x)\}\subseteq A}[thm],\FormulaRefAuto{\{x\in A\mid P(x)\}\subseteq A}[thm]}}
 \proofstepwidestar[2]{t^{-1}[I'_{\psi(p)}]\subseteq T}{\FormulaRefAuto{F\colon A\to B\dsep C\subseteq B\vdash F^{-1}[C]\subseteq A}[thm]{2,\rAEb{24}}}
 \proofstepwidestar[1,2,3]{t^{-1}[\{\psi(p)\}]\subseteq T}{\FormulaRefAuto{F\colon A\to B\dsep C\subseteq B\vdash F^{-1}[C]\subseteq A}[thm]{2,\FormulaRefAuto{a\in A\vdash\{a\}\subseteq A}[thm]{7}}}
 \proofstepwidestar[2,8]{x\in t^{-1}[I'_{\psi(p)}]\leftrightarrow t(x)\preceq\psi(p)}{\rChain{\FormulaRefAuto{BoundedPreimageMembership}[thm]{2,\rAEb{24},8},\FormulaRefAuto{BoundedSeparationMembership}[thm]{9}}}
 \proofstepwidestar[1,2,8]{v(x)\in I_p\leftrightarrow v(x)\leq p}{\FormulaRefAuto{BoundedSeparationMembership}[thm]{10}}
 \proofstepwidestar[1,2,3,4,8]{x\in t^{-1}[I'_{\psi(p)}]\leftrightarrow v(x)\in I_p}{\rChain{27,\FormulaRefAuto{P\leftrightarrow Q\vdash Q\leftrightarrow P}[thm]{15},\FormulaRefAuto{P\leftrightarrow Q\vdash Q\leftrightarrow P}[thm]{28}}}
 \proofstepwidestar[1,2,3,4]{\forall x\in T\,(x\in t^{-1}[I'_{\psi(p)}]\leftrightarrow v(x)\in I_p)}{\rUI{\rRI{8,29}}}
 \proofstepwidestar[1,2,3,4]{v^{-1}[I_p]=t^{-1}[I'_{\psi(p)}]}{\FormulaRefAuto{PreimageEqFromPointwiseIff}[thm]{6,25,\rAEa{24},30}}
 \proofstepwidestar[1,2,3,8]{x\in t^{-1}[\{\psi(p)\}]\leftrightarrow t(x)=\psi(p)}{\rChain{\FormulaRefAuto{BoundedPreimageMembership}[thm]{2,\FormulaRefAuto{a\in A\vdash\{a\}\subseteq A}[thm]{7},8},\FormulaRefAuto{x\in\{a\}\eqvdash x=a}[thm]}}
 \proofstepwidestar[]{v(x)\in\{p\}\leftrightarrow v(x)=p}{\FormulaRefAuto{x\in\{a\}\eqvdash x=a}[thm]}
 \proofstepwidestar[1,2,3,8]{x\in t^{-1}[\{\psi(p)\}]\leftrightarrow v(x)\in\{p\}}{\rChain{32,\FormulaRefAuto{P\leftrightarrow Q\vdash Q\leftrightarrow P}[thm]{23},\FormulaRefAuto{P\leftrightarrow Q\vdash Q\leftrightarrow P}[thm]{33}}}
 \proofstepwidestar[1,2,3]{\forall x\in T\,(x\in t^{-1}[\{\psi(p)\}]\leftrightarrow v(x)\in\{p\})}{\rUI{\rRI{8,34}}}
 \proofstepwidestar[3]{\{p\}\subseteq Q}{\FormulaRefAuto{a\in A\vdash\{a\}\subseteq A}[thm]{3}}
 \proofstepwidestar[1,2,3]{v^{-1}[\{p\}]=t^{-1}[\{\psi(p)\}]}{\FormulaRefAuto{PreimageEqFromPointwiseIff}[thm]{6,26,36,35}}
 \proofstepwidestar[1,2,3,4]{\begin{aligned}[t]&v^{-1}[I_p]=t^{-1}[I'_{\psi(p)}]\land{}\\[-2pt]& v^{-1}[\{p\}]=t^{-1}[\{\psi(p)\}]\end{aligned}}{\rAI{31,37}}
 \closeproofpartwide
\end{tabproofsplitwide}

\Needspace{14\baselineskip}
\FormulaThmDeltaK[Fasergrößen aus korrespondierenden unteren Mengen]{%
 \begin{aligned}[t]
 &\PartOrd{Q,\leq}\dsep\PartOrd{R,\preceq}\dsep\Finite Q\dsep\Finite S\dsep\Finite T\dsep{}\\[-2pt]
 &r\colon S\to Q\dsep t\colon T\to R\dsep\psi\colon Q\bij R\dsep{}\\[-2pt]
 &\forall p,q\in Q\,(p\leq q\leftrightarrow\psi(p)\preceq\psi(q))\dsep{}\\[-2pt]
 &\forall p\in Q\,\FiniteCard(r^{-1}[I_p])=\FiniteCard(t^{-1}[I'_{\psi(p)}])\vdash{}\\[-2pt]
 &\forall p\in Q\,\exists u\,(u\colon r^{-1}[\{p\}]\bij t^{-1}[\{\psi(p)\}])
 \end{aligned}
}{FinitePosetCumulativeTransportBijection}{%
 \DeltaRow{Mengen}{Q\dsep R\dsep S\dsep T\dsep p\dsep q}
 \DeltaRow{Relationsoperatoren}{\leq\dsep\preceq}
 \DeltaRow{Funktionen}{r\dsep t\dsep\psi\dsep u\dsep v=\psi^{-1}\circ t}
 \DeltaRow{Abkürzungen}{I_p=\{q\in Q\mid q\leq p\}\dsep I'_q=\{b\in R\mid b\preceq q\}}
}
\begin{tabproofwide}
 \proofstepwidestar[1]{\PartOrd{Q,\leq}\land\PartOrd{R,\preceq}}{\rA}
 \proofstepwidestar[2]{\Finite Q\land\Finite S\land\Finite T}{\rA}
 \proofstepwidestar[3]{r\colon S\to Q\land t\colon T\to R\land\psi\colon Q\bij R}{\rA}
 \proofstepwidestar[4]{\forall p,q\in Q\,(p\leq q\leftrightarrow\psi(p)\preceq\psi(q))}{\rA}
 \proofstepwidestar[5]{\forall p\in Q\,\FiniteCard(r^{-1}[I_p])=\FiniteCard(t^{-1}[I'_{\psi(p)}])}{\rA}
 \proofstepwidestar[3]{\psi^{-1}\colon R\to Q}{\FormulaRefAuto{InverseFunctionType}[thm]{\rAEb{\rAEb{3}}}}
 \proofstepwidestar[3]{v\colon T\to Q}{\FormulaRefAuto{CompositionDef}[def]{\rAEa{\rAEb{3}},6}}
 \proofstepwidestar[8]{p\in Q}{\rA}
 \proofstepwidestar[3,4,8]{\begin{aligned}[t]&v^{-1}[I_p]=t^{-1}[I'_{\psi(p)}]\land{}\\[-2pt]& v^{-1}[\{p\}]=t^{-1}[\{\psi(p)\}]\end{aligned}}{\FormulaRefAuto{OrderPullbackLowerFibers}[thm]{\rAEb{\rAEb{3}},\rAEa{\rAEb{3}},8,4}}
 \proofstepwidestar[5,8]{\FiniteCard(r^{-1}[I_p])=\FiniteCard(t^{-1}[I'_{\psi(p)}])}{\rRE{\rUE{5},8}}
 \proofstepwidestar[3,4,5,8]{\FiniteCard(r^{-1}[I_p])=\FiniteCard(v^{-1}[I_p])}{\rIE{\FormulaRefAuto{a=b\vdash b=a}[thm]{\rAEa{9}},10}}
 \proofstepwidestar[3,4,5]{\forall p\in Q\,\FiniteCard(r^{-1}[I_p])=\FiniteCard(v^{-1}[I_p])}{\rUI{\rRI{8,11}}}
 \proofstepwidestar[1,2,3,4,5]{\forall p\in Q\,\exists u\,(u\colon r^{-1}[\{p\}]\bij v^{-1}[\{p\}])}{\FormulaRefAuto{FinitePosetCumulativeFiberBijection}[thm]{\rAEa{1},\rAEa{2},\rAEa{\rAEb{2}},\rAEb{\rAEb{2}},\rAEa{3},7,12}}
 \proofstepwidestar[1,2,3,4,5,8]{\exists u\,(u\colon r^{-1}[\{p\}]\bij v^{-1}[\{p\}])}{\rRE{\rUE{13},8}}
 \proofstepwidestar[1,2,3,4,5,8]{\exists u\,(u\colon r^{-1}[\{p\}]\bij t^{-1}[\{\psi(p)\}])}{\rIE{\rAEb{9},14}}
 \proofstepwidestar[1,2,3,4,5]{\forall p\in Q\,\exists u\,(u\colon r^{-1}[\{p\}]\bij t^{-1}[\{\psi(p)\}])}{\rUI{\rRI{8,15}}}
\end{tabproofwide}

